\documentclass[9pt]{extarticle}

\usepackage[utf8]{inputenc}
\usepackage[a4paper, margin=1.2cm]{geometry}
\usepackage[colorlinks=true, urlcolor=blue]{hyperref}
\usepackage{enumitem}
\usepackage{changepage}

\pagestyle{empty}

\setlist[itemize]{label=\textbullet, noitemsep, topsep=0pt}
\setlength{\parindent}{0pt}

% actual document shit

\begin{document}

\begin{center}
	{\huge\bfseries Paul Parisot}
	\\[9pt]
	{\Large Ingénieur Logiciel - Programmation Système - Développeur Backend}
	\\
\end{center}

% section for contact

\par\vspace{-5pt}
\rule{\linewidth}{0.4pt}
\par\vspace{-5pt}

\begin{center}
	Github: \href{https://github.com/paulogarithm}{paulogarithm}
	\\
	\href{mailto:paul.parisot@epitech.eu}{paul.parisot@epitech.eu}
	\\
\end{center}

\par\vspace{-10pt}
\rule{\linewidth}{0.4pt}
\par\vspace{7pt}

% section for the about us

{\large\bfseries À propos}
\par\vspace{-8pt}
\rule{\linewidth}{0.4pt}
\par\vspace{5pt}

Ingénieur logiciel orienté systèmes, avec de solides bases en C/C++, Go et langages bas niveau.
\\
Expérience dans la conception de langages de programmation, moteurs de jeux et systèmes distribués.
\\
Contributeur open-source avec plus de 10 000 lignes de code sur des packages critiques.
\\
Actuellement en préparation de l'agrégation d'informatique, concours national d'enseignement très sélectif, afin d'approfondir mes bases théoriques en algorithmique des graphes, théorie de la complexité et automates.
\par\vspace{5pt}

% section for skills

{\large\bfseries Compétences}
\par\vspace{-8pt}
\rule{\linewidth}{0.4pt}
\par\vspace{5pt}

\hspace*{5pt}{\bfseries Techniques}: C/C++, Golang, Python, Flutter, JavaScript/TypeScript, Haskell, Lua(u), OCaml, Assembleur
\\
\hspace*{5pt}{\bfseries Autres}: Outils Git, DevOps (Docker, Kubernetes, ...), CI/CD, Environnement Linux

\par\vspace{5pt}

% section for the professional experience

{\large\bfseries Expérience Professionnelle}
\par\vspace{-8pt}
\rule{\linewidth}{0.4pt}
\par\vspace{5pt}

{\bfseries Gno.land - Ingénieur Logiciel Junior (Open-Source)}
\begin{itemize}[noitemsep, topsep=0pt]
	\item Développement de 6+ packages core en Gnolang (variante de Go), incluant printf et une compatibilité avec l'EVM.
	\item Détection et correction de {\bfseries 6 bugs critiques} causant des crashs runtime ; {\bfseries ~10k lignes de code} contribuées.
	\item Rédaction de tests exhaustifs sous des politiques strictes de CI/CD et de revue de code.
\end{itemize}

\par\vspace{5pt}

{\bfseries Faurecia / Forvia - Ingénieur Démo Logicielle}
\begin{itemize}[noitemsep, topsep=0pt]
	\item Développement d'une {\bfseries application démo avec Unity}, présentée au {\bfseries CES 2024}, simulant navigation GPS \& UI.
	\item Conception d'outils interactifs comme un {\bfseries moteur de calcul d'itinéraire GPS} et une interface de recherche de POI.
	\item Livraison de fonctionnalités prêtes pour la production {\bfseries en avance sur le planning}, salué pour la rapidité et la qualité du code.
\end{itemize}

\par\vspace{5pt}

{\bfseries Epitech / IONIS - Assistant Technique \& Formateur}
\begin{itemize}[noitemsep, topsep=0pt]
	\item Enseignement à {\bfseries plus de 400 étudiants} en C, C++, asm et Haskell ; correction et évaluation des rendus.
	\item Création de {\bfseries contenus pédagogiques interactifs} (quiz, cours inversés, tutoriels).
	\item Application de {\bfseries méthodes pédagogiques} (constructivisme, cognitivisme…) pour renforcer l'engagement des étudiants.
	\item Contribution à un {\bfseries taux de réussite historiquement élevé} au sein des promotions coachées par l'ASTEK.
\end{itemize}

\par\vspace{5pt}

% personal projects

{\large\bfseries Projets Personnels}
\par\vspace{-8pt}
\rule{\linewidth}{0.4pt}
\par\vspace{5pt}

\underline{\bfseries Sledge -- Moteur de migration, synchronisation et orchestration multi-cloud}
\par
\begin{adjustwidth}{5pt}{0pt}
    Développement et finalisation d'un outil générique de migration multi-cloud gérant des structures de données complexes et des interactions avec des bases de données. \\
    Architecture et livraison d'une plateforme backend complète en Go, optimisée pour des opérations de flux de données à haut volume, via une couche d'abstraction cloud personnalisée ({\ttfamily provicol}) supportant AWS S3 et GCP. \\
    Utilisation de l'API Docker à l'exécution pour créer et orchestrer dynamiquement des conteneurs d'environnements concurrents. \\
    Python3 a été grandement utilisé pour la cli en interne, et pour les jeux de tests de Sledge. \\
\end{adjustwidth}
\par
{\bfseries Compétences :} Go, API Docker, Python, SQL, intégrité des données, cloud computing, automatisation système

\par\vspace{8pt}

\underline{\bfseries \href{https://github.com/paulogarithm/zzzcript}{ZZZcript} -- Langage fonctionnel statiquement typé}
\par
\begin{adjustwidth}{5pt}{0pt}
	Conception d'un {\bfseries lexer, parser et interpréteur AST} from scratch. \\
	Écriture d'un moteur d'évaluation récursif en {\bfseries Go}, inspiré de Haskell et Golang.
\end{adjustwidth}
\par
{\bfseries Compétences :} Go, Haskell, théorie du parsing, conception d'AST, sémantique fonctionnelle, architecture d'interpréteur

\par\vspace{8pt}

\underline{\bfseries \href{https://github.com/paulogarithm/cclass}{cclass} -- Moteur de mutation mémoire à l'exécution}
\par
\begin{adjustwidth}{5pt}{0pt}
    Conception d'un framework en C pur modifiant dynamiquement des instructions assembleur compilées directement en mémoire. \\
    Utilisation d'appels système {\bfseries \texttt{mmap}} avec des flags de protection précis pour contourner le surcoût du dispatch par vtable. \\
    Application de règles d'alignement mémoire strictes pour garantir une stabilité absolue sur les cibles d'exécution matérielles.
\end{adjustwidth}
\par
{\bfseries Compétences :} C, systèmes Linux, assembleur x86, mmap, sûreté mémoire, débogage bas niveau

\par\vspace{8pt}

\underline{\bfseries \href{https://github.com/paulogarithm/blorox}{rodox} -- Moteur de jeu 3D modulaire inspiré de Roblox}
\par
\begin{adjustwidth}{5pt}{0pt}
	Conception d'un moteur de jeu en C avec modules de rendu et de logique chargés à l'exécution. \\
	Intégration de scripting Lua et imitation de l'API Roblox pour la compatibilité. \\
	En développement actif, finalisation prévue à l'arrivée des fonctionnalités de reflection de C++26.
\end{adjustwidth}
\par
{\bfseries Compétences :} C, C++, Lua, architecture ECS, systèmes de plugins, liaison à l'exécution, conception de moteur

\par\vspace{8pt}

% education

{\large\bfseries Formation}
\par\vspace{-8pt}
\rule{\linewidth}{0.4pt}
\par\vspace{5pt}

{\bfseries Epitech, Paris} \hfill 2022 -- 2027 \\
Diplôme d'ingénieur en technologies de l'information (diplôme de niveau Master) \\
Certification Niveau 7 - Expert en ingénierie logicielle

\vspace{6pt}

{\bfseries McGill University, Montréal} \hfill 2025 -- 2026 \\
Certificat en Management \\
Cours suivis : Comptabilité, Administration des affaires, Management organisationnel

\end{document}
